\phantomsection
\section*{Глава 2. Проектирование веб-приложения}
\addcontentsline{toc}{section}{Глава 2. Проектирование веб-приложения}
\stepcounter{section}

\subsection{Анализ целевой аудитории и ролевая модель}

Разрабатываемое приложение предназначено для трёх категорий пользователей языкового
центра, каждая из которых имеет собственные цели и модель использования системы.

\textbf{Студенты} - основные потребители образовательных услуг. Им необходимо
просматривать актуальное расписание занятий, самостоятельно записываться на курсы
с учётом доступных мест, отслеживать собственную успеваемость и посещаемость,
а также получать учебные материалы. Для данной группы важны простота навигации
и наличие личного кабинета с агрегированными показателями прогресса.

\textbf{Преподаватели} - ведут занятия и несут ответственность за оценивание
студентов. Нуждаются в удобном интерфейсе для просмотра списка записавшихся
на конкретное занятие, выставления оценок и отметок о посещаемости, публикации
учебных материалов и экспорта журнала успеваемости.

\textbf{Администраторы} - управляют учебным процессом и структурой данных.
Отвечают за создание и редактирование курсов, формирование расписания занятий,
управление учётными записями пользователей, назначение ролей и рассмотрение заявок
студентов на получение роли преподавателя.

Исходя из описания аудитории, в системе определены четыре уровня доступа.

\textbf{Гость} - неаутентифицированный посетитель. Может просматривать каталог
курсов и общедоступную информацию, но не имеет доступа к расписанию, записи
на занятия и личному кабинету.

\textbf{Студент (Student)} - зарегистрированный участник с базовым уровнем
доступа. Может просматривать расписание, записываться на занятия и отменять запись,
оставлять оценки и комментарии к курсам, просматривать учебные материалы и личный
прогресс, а также подавать заявку на роль преподавателя.

\textbf{Преподаватель (Teacher)} - роль, назначаемая администратором или
присваиваемая после одобрения заявки. Обладает всеми правами студента,
а также может выставлять оценки и отметки посещаемости, публиковать и редактировать
учебные материалы, экспортировать журнал успеваемости в формате CSV.

\textbf{Администратор (Admin)} - привилегированная роль с полным доступом
ко всем разделам системы. Управляет курсами, занятиями и пользователями,
назначает роли, рассматривает заявки на роль преподавателя.

Матрица прав доступа по основным операциям системы приведена
в таблице~\ref{tab:roles}.

\begin{table}[ht!]
    \centering
    \caption{Матрица прав доступа по ролям}
    \label{tab:roles}
    \small
    \begin{tabularx}{\textwidth}{|X|c|c|c|c|}
        \hline
        \textbf{Операция} &
        \textbf{Гость} &
        \textbf{Студент} &
        \textbf{Препод.} &
        \textbf{Администр.} \\ \hline
        Просмотр каталога курсов                & Да  & Да  & Да  & Да  \\ \hline
        Регистрация / вход в систему            & Да  & --  & --  & --  \\ \hline
        Просмотр расписания занятий             & --  & Да  & Да  & Да  \\ \hline
        Запись на занятие / отмена записи       & --  & Да  & --  & --  \\ \hline
        Просмотр учебных материалов             & --  & Да  & Да  & Да  \\ \hline
        Оценка курса и комментирование          & --  & Да  & Да  & Да  \\ \hline
        Просмотр личного прогресса              & --  & Да  & --  & --  \\ \hline
        Подача заявки на роль преподавателя     & --  & Да  & --  & --  \\ \hline
        Выставление оценок и посещаемости       & --  & --  & Да  & Да  \\ \hline
        Публикация и редактирование материалов  & --  & --  & Да  & Да  \\ \hline
        Экспорт журнала успеваемости (CSV)      & --  & --  & Да  & Да  \\ \hline
        Управление курсами (CRUD)               & --  & --  & --  & Да  \\ \hline
        Управление занятиями (CRUD)             & --  & --  & --  & Да  \\ \hline
        Управление пользователями и ролями      & --  & --  & --  & Да  \\ \hline
        Рассмотрение заявок на роль преподавателя & -- & --  & --  & Да  \\ \hline
    \end{tabularx}
\end{table}

\subsection{Функциональные требования}

На основе ролевой модели выделены ключевые варианты использования системы,
представленные в виде диаграммы на рисунке~\ref{fig:use_case}.

\begin{figure}[ht!]
    \centering
    \includegraphics[width=0.65\textwidth]{use_case_diagram}
    \caption{Диаграмма вариантов использования LangCenter}
    \label{fig:use_case}
\end{figure}

Диаграмма фиксирует взаимодействие четырёх акторов со следующими группами
вариантов использования:

\begin{itemize}
    \item \textbf{каталог и регистрация} - просмотр курсов и страниц курса
          доступен всем; регистрация и вход - только гостям;
    \item \textbf{запись и расписание} - просмотр расписания, запись
          на занятие и её отмена (студент и выше);
    \item \textbf{учебные материалы} - просмотр и скачивание (студент),
          публикация, редактирование и удаление (преподаватель и выше);
    \item \textbf{успеваемость} - просмотр личного прогресса (студент),
          выставление оценок и экспорт журнала в CSV (преподаватель и выше);
    \item \textbf{администрирование} - полный CRUD курсов, занятий,
          управление ролями и рассмотрение заявок (только администратор).
\end{itemize}

Перечень URL-маршрутов приложения, соответствующих вариантам использования,
приведён в таблице~\ref{tab:routes}.

\begin{table}[ht!]
    \centering
    \caption{Маршруты веб-приложения LangCenter}
    \label{tab:routes}
    \small
    \begin{tabularx}{\textwidth}{|p{1.5cm}|p{5.2cm}|p{2.4cm}|X|}
        \hline
        \textbf{Метод} & \textbf{URL} & \textbf{Мин. роль} & \textbf{Назначение} \\ \hline
        GET            & /                                          & Гость       & Главная страница \\ \hline
        GET/ POST       & /auth/register                             & Гость       & Регистрация \\ \hline
        GET/ POST       & /auth/login                                & Гость       & Вход в систему \\ \hline
        GET            & /auth/logout                               & Студент     & Выход \\ \hline
        GET            & /courses                                   & Гость       & Каталог курсов с поиском и фильтрами \\ \hline
        GET            & /courses/\textlangle{}id\textrangle{}      & Гость       & Страница курса: материалы, оценки, комментарии \\ \hline
        GET            & /schedule                                  & Студент     & Недельное расписание с фильтрами \\ \hline
        POST           & /lessons/\textlangle{}id\textrangle{}/enroll   & Студент & Запись на занятие \\ \hline
        POST           & /lessons/\textlangle{}id\textrangle{}/unenroll & Студент & Отмена записи \\ \hline
        GET            & /my-lessons                                & Студент     & Список собственных занятий \\ \hline
        GET            & /my-progress                               & Студент     & Личная статистика успеваемости \\ \hline
        GET/ POST       & /request-teacher                           & Студент     & Заявка на роль преподавателя \\ \hline
        GET            & /teacher/                                  & Преподав.   & Личный кабинет преподавателя \\ \hline
        GET/ POST       & /teacher/lessons/\textlangle{}id\textrangle{}/grades & Преподав. & Выставление оценок и посещаемости \\ \hline
        GET            & /teacher/lessons/\textlangle{}id\textrangle{}/grades /export & Преподав. & Экспорт журнала в CSV \\ \hline
        GET            & /admin/                                    & Администр.  & Панель администратора (статистика) \\ \hline
        GET/ POST       & /admin/courses/new                         & Администр.  & Создание курса \\ \hline
        GET/ POST       & /admin/lessons/new                         & Администр.  & Создание занятия \\ \hline
        POST           & /admin/users/\textlangle{}id\textrangle{}/role & Администр. & Смена роли пользователя \\ \hline
        POST           & /admin/requests/\textlangle{}id\textrangle{}/approve & Администр. & Одобрение заявки на преподавателя \\ \hline
        GET            & /notifications/                            & Студент     & Список уведомлений \\ \hline
        GET            & /notifications/api/poll                    & Студент     & Polling новых уведомлений (JSON) \\ \hline
    \end{tabularx}
\end{table}

\subsection{Проектирование модели данных}

База данных приложения построена на реляционной модели и включает десять сущностей.
ER-диаграмма, отражающая состав полей и связи между таблицами, представлена
на рисунке~\ref{fig:er_diagram}.

\begin{figure}[ht!]
    \centering
    \includegraphics[width=\textwidth]{ER}
    \caption{ER-диаграмма базы данных LangCenter}
    \label{fig:er_diagram}
\end{figure}

Связи между основными сущностями имеют следующий характер:

\begin{itemize}
    \item \texttt{users} $\leftrightarrow$ \texttt{courses}: один преподаватель
          ведёт несколько курсов (1:N, внешний ключ \texttt{teacher\_id});
    \item \texttt{courses} $\leftrightarrow$ \texttt{lessons}: один курс включает
          несколько занятий (1:N, внешний ключ \texttt{course\_id});
    \item \texttt{lessons} $\leftrightarrow$ \texttt{enrollments}: на одно занятие
          записываются несколько студентов (1:N, внешний ключ \texttt{lesson\_id});
    \item \texttt{users} $\leftrightarrow$ \texttt{enrollments}: один студент
          может быть записан на несколько занятий; пара (\texttt{student\_id},
          \texttt{lesson\_id}) уникальна;
    \item \texttt{enrollments} $\leftrightarrow$ \texttt{grades}: каждая запись
          может иметь одну оценку (1:1, внешние ключи \texttt{student\_id}
          и \texttt{lesson\_id}).
\end{itemize}

Структура таблицы \texttt{users} представлена в таблице~\ref{tab:db_users}.

{\small
\begin{longtable}{|p{3.8cm}|p{2.8cm}|p{3.5cm}|p{4.4cm}|}
    \caption{Структура таблицы пользователей (\texttt{users})}
    \label{tab:db_users} \\
    \hline
    \textbf{Поле} & \textbf{Тип} & \textbf{Ограничение} & \textbf{Описание} \\ \hline
    \endfirsthead
    \tabcont{\ref{tab:db_users}}
    \textbf{Поле} & \textbf{Тип} & \textbf{Ограничение} & \textbf{Описание} \\ \hline
    \endhead
    \hline \endfoot
    \hline \endlastfoot
    id            & INTEGER      & PRIMARY KEY, SERIAL        & Уникальный идентификатор \\ \hline
    first\_name   & VARCHAR(64)  & NOT NULL                   & Имя пользователя \\ \hline
    last\_name    & VARCHAR(64)  & NOT NULL                   & Фамилия пользователя \\ \hline
    email         & VARCHAR(120) & UNIQUE, NOT NULL           & Адрес электронной почты \\ \hline
    password\_hash & VARCHAR(256) & NOT NULL                  & Хеш пароля (werkzeug pbkdf2) \\ \hline
    role          & VARCHAR(20)  & NOT NULL, DEFAULT 'student' & Роль: admin, teacher, student \\ \hline
    created\_at   & TIMESTAMP    & DEFAULT NOW()              & Дата и время регистрации \\ \hline
\end{longtable}
}

Структура таблицы \texttt{courses} представлена в таблице~\ref{tab:db_courses}.

{\small
\begin{longtable}{|p{3.8cm}|p{2.8cm}|p{3.5cm}|p{4.4cm}|}
    \caption{Структура таблицы курсов (\texttt{courses})}
    \label{tab:db_courses} \\
    \hline
    \textbf{Поле} & \textbf{Тип} & \textbf{Ограничение} & \textbf{Описание} \\ \hline
    \endfirsthead
    \tabcont{\ref{tab:db_courses}}
    \textbf{Поле} & \textbf{Тип} & \textbf{Ограничение} & \textbf{Описание} \\ \hline
    \endhead
    \hline \endfoot
    \hline \endlastfoot
    id          & INTEGER      & PRIMARY KEY, SERIAL         & Уникальный идентификатор курса \\ \hline
    name        & VARCHAR(200) & NOT NULL                    & Название курса \\ \hline
    language    & VARCHAR(64)  & NOT NULL                    & Изучаемый язык \\ \hline
    level       & VARCHAR(10)  & NOT NULL (A1--C2)           & Уровень владения языком \\ \hline
    description & TEXT         & NULL                        & Описание курса \\ \hline
    teacher\_id & INTEGER      & FK → users.id, NOT NULL     & Ссылка на преподавателя \\ \hline
    is\_active  & BOOLEAN      & DEFAULT TRUE                & Флаг активности курса \\ \hline
    created\_at & TIMESTAMP    & DEFAULT NOW()               & Дата создания \\ \hline
\end{longtable}
}

Структура таблицы \texttt{lessons} представлена в таблице~\ref{tab:db_lessons}.

\begin{table}[ht!]
    \centering
    \caption{Структура таблицы занятий (\texttt{lessons})}
    \label{tab:db_lessons}
    \small
    \begin{tabularx}{\textwidth}{|p{3.0cm}|p{2.5cm}|p{3.5cm}|X|}
        \hline
        \textbf{Поле} & \textbf{Тип} & \textbf{Ограничение} & \textbf{Описание} \\ \hline
        id           & INTEGER   & PRIMARY KEY, SERIAL      & Уникальный идентификатор \\ \hline
        course\_id   & INTEGER   & FK → courses.id, NOT NULL & Ссылка на курс \\ \hline
        teacher\_id  & INTEGER   & FK → users.id, NOT NULL  & Ссылка на преподавателя \\ \hline
        starts\_at   & TIMESTAMP & NOT NULL                 & Дата и время начала занятия \\ \hline
        duration\_min & INTEGER  & NOT NULL, DEFAULT 60     & Продолжительность, мин \\ \hline
        room         & VARCHAR(32) & NULL                   & Номер аудитории \\ \hline
        capacity     & INTEGER   & NOT NULL, DEFAULT 10     & Максимальное число мест \\ \hline
    \end{tabularx}
\end{table}

Структура таблицы \texttt{enrollments} представлена в таблице~\ref{tab:db_enrollments}.

\begin{table}[ht!]
    \centering
    \caption{Структура таблицы записей (\texttt{enrollments})}
    \label{tab:db_enrollments}
    \small
    \begin{tabularx}{\textwidth}{|p{3.0cm}|p{2.5cm}|p{3.5cm}|X|}
        \hline
        \textbf{Поле} & \textbf{Тип} & \textbf{Ограничение} & \textbf{Описание} \\ \hline
        id          & INTEGER   & PRIMARY KEY, SERIAL                          & Уникальный идентификатор \\ \hline
        student\_id & INTEGER   & FK → users.id, NOT NULL                      & Ссылка на студента \\ \hline
        lesson\_id  & INTEGER   & FK → lessons.id, NOT NULL                    & Ссылка на занятие \\ \hline
        status      & VARCHAR   & DEFAULT 'active'                             & Статус: active, cancelled \\ \hline
        enrolled\_at & TIMESTAMP & DEFAULT NOW()                               & Дата и время записи \\ \hline
        \multicolumn{4}{|l|}{UNIQUE (student\_id, lesson\_id)} \\ \hline
    \end{tabularx}
\end{table}

Структура таблицы \texttt{grades} представлена в таблице~\ref{tab:db_grades}.

\begin{table}[ht!]
    \centering
    \caption{Структура таблицы оценок (\texttt{grades})}
    \label{tab:db_grades}
    \small
    \begin{tabularx}{\textwidth}{|p{3.0cm}|p{2.5cm}|p{3.5cm}|X|}
        \hline
        \textbf{Поле} & \textbf{Тип} & \textbf{Ограничение} & \textbf{Описание} \\ \hline
        id          & INTEGER   & PRIMARY KEY, SERIAL       & Уникальный идентификатор \\ \hline
        student\_id & INTEGER   & FK → users.id, NOT NULL   & Ссылка на студента \\ \hline
        lesson\_id  & INTEGER   & FK → lessons.id, NOT NULL & Ссылка на занятие \\ \hline
        score       & FLOAT     & NULL, CHECK (0--100)      & Балл (0--100) \\ \hline
        attended    & BOOLEAN   & DEFAULT FALSE             & Отметка о посещении \\ \hline
        notes       & TEXT      & NULL                      & Комментарий преподавателя \\ \hline
        created\_at & TIMESTAMP & DEFAULT NOW()             & Дата выставления \\ \hline
        \multicolumn{4}{|l|}{UNIQUE (student\_id, lesson\_id)} \\ \hline
    \end{tabularx}
\end{table}

Помимо перечисленных, база данных включает вспомогательные таблицы:
\texttt{course\_materials} (учебные материалы с возможностью файлового вложения),
\texttt{course\_ratings} (оценки курса по шкале 1--5, уникальные для пары
пользователь/курс), \texttt{course\_comments} (комментарии к курсам),
\texttt{notifications} (внутренние уведомления с флагом прочтения) и
\texttt{teacher\_requests} (заявки на роль преподавателя со статусами
\texttt{pending}, \texttt{approved}, \texttt{rejected}).

Средний балл студента по курсу вычисляется как среднее арифметическое
всех выставленных оценок:

\begin{equation}
    \bar{P} = \frac{1}{N} \sum_{i=1}^{N} s_i,
    \label{eq:avg_grade}
\end{equation}

\begin{where}
    \item $s_i$~--- балл за $i$-е занятие (от 0 до 100);
    \item $N$~--- количество занятий, за которые выставлена оценка.
\end{where}

\subsection{Архитектура приложения}

Приложение построено по трёхуровневой архитектуре
(рисунок~\ref{fig:architecture}): клиентский уровень (браузер), серверный
уровень (Flask-приложение) и уровень хранения данных (PostgreSQL и файловая
система для учебных материалов).

\begin{figure}[ht!]
    \centering
    \includegraphics[width=0.88\textwidth]{app_arch}
    \caption{Архитектура приложения LangCenter}
    \label{fig:architecture}
\end{figure}

Серверный уровень организован по шаблону Application Factory:
функция \texttt{create\_app()} создаёт и конфигурирует экземпляр приложения,
подключает расширения (SQLAlchemy, Migrate, LoginManager, CSRFProtect)
и регистрирует блупринты. Маршруты сгруппированы в шесть блупринтов:
main (каталог, расписание, запись, прогресс), auth
(аутентификация, профиль), admin (панель администратора),
teacher (оценивание, экспорт), notifications (уведомления)
и соответствующие шаблоны Jinja2.

Структура каталога проекта показана в листинге~\ref{lst:structure}.

\begin{lstlisting}[caption={Структура проекта LangCenter}, label={lst:structure}]
mirs/
  App/
    app/
      __init__.py        # Application Factory
      extensions.py      # db, migrate, login_manager, csrf
      config.py          # DevelopmentConfig / ProductionConfig
      commands.py        # flask seed (тестовые данные)
      models/
        user.py          # User, Role enum
        course.py        # Course, LanguageLevel enum
        lesson.py        # Lesson
        enrollment.py    # Enrollment, EnrollmentStatus
        grade.py         # Grade
        course_material.py  # CourseMaterial (файлы)
        course_rating.py    # CourseRating
        course_comment.py   # CourseComment
        notification.py     # Notification
        teacher_request.py  # TeacherRequest
      main/
        routes.py        # Каталог, расписание, запись, прогресс
        forms.py         # MaterialForm, TeacherRequestForm
      auth/
        routes.py        # register, login, logout, profile
        forms.py         # LoginForm, RegisterForm
        decorators.py    # admin_required, teacher_required
      admin/
        routes.py        # CRUD курсов, занятий, пользователей
        forms.py         # CourseForm, LessonForm
      teacher/
        routes.py        # grade_lesson, export_grades
      notifications/
        routes.py        # index, poll, mark_read
        utils.py         # notify()
      templates/
        base.html        # Базовый шаблон (navbar, toast)
        ...
      static/
        css/style.css    # Кастомные стили (sketchy-дизайн)
    migrations/          # Alembic-миграции
    Dockerfile
    requirements.txt
  docker-compose.yml
  .env
\end{lstlisting}

\subsection{Проектирование интерфейса}

Интерфейс приложения проектировался из расчёта трёх ключевых сценариев:
просмотр каталога и запись на занятие, работа преподавателя с оценками,
управление структурой центра администратором.

\textbf{Каталог курсов} - основная точка входа. Страница содержит строку
текстового поиска по названию и описанию, фильтры по языку и уровню владения,
сетку карточек курсов с названием, уровнем, именем преподавателя и средним
рейтингом. При наличии активного поиска отображается количество найденных
результатов.

\textbf{Страница курса} - детальная информация: описание, список учебных
материалов (с возможностью скачивания файлов), ближайшие занятия с кнопкой
записи, блок звёздного рейтинга и секция комментариев. Для преподавателя
и администратора дополнительно отображаются кнопки управления материалами.

\textbf{Расписание} - недельный вид с навигацией по неделям и фильтрами
по языку, уровню и преподавателю. Занятия отображаются в виде карточек
с указанием времени, аудитории, количества доступных мест и кнопкой записи.

\textbf{Журнал оценок} - доступен преподавателю. Таблица со списком
записавшихся студентов, полями для ввода балла (0--100) и чекбоксом
посещаемости. Кнопка экспорта в CSV расположена в верхнем правом углу.

Макеты двух ключевых страниц - каталога курсов и расписания занятий -
представлены на рисунках~\ref{fig:mockup_catalog} и~\ref{fig:mockup_schedule}
соответственно.

\begin{figure}[ht!]
    \centering
    \includegraphics[width=0.82\textwidth]{catalog}
    \caption{Макет страницы каталога курсов}
    \label{fig:mockup_catalog}
\end{figure}

\begin{figure}[ht!]
    \centering
    \includegraphics[width=0.82\textwidth]{calenfar}
    \caption{Макет страницы расписания занятий}
    \label{fig:mockup_schedule}
\end{figure}
